\lecture{7}{Wed 14 Sep 2022 11:31}{Closed Sets}

\begin{example}
	(countable union of open sets need not be open)
	If $U_n = B_{\frac{1}{n}}(x), x \in \mathbb{R}^p$. Then \[
		\bigcap_{n = 1}^\infty  B_{\frac{1}{n}}(x) = \{x\} 
	.\] 
	\begin{proof}
		If $y \neq x$ then $\|y-x\| >0$. Exists $n \in \mathbb{N}: 0 < \frac{1}{n} < \|y-x\|$. Hence $y \not\in B_{\frac{1}{n}}(x)$.
	\end{proof}
\end{example}

\begin{proposition}
	A set $E \subset \mathbb{R}^p$ is open if and only if $E$ is a union of open balls.
\end{proposition}
\begin{proof}
	($\impliedby$) is clear by the previous proposition, since open balls are open.

	Now we prove ($\implies$ ). For any $x \in E$, there exists $r = r_x > 0$ such that \[
		G_x = B_{r_x}(x) \subset E
	.\] 
	Then $E = \bigcup_{x \in E} G_x$. Indeed, since $G_x \subset E$, the union is $\subset E$. For any $x_0 \in E$, we have $x_0 \in G_{x_0} \subset \bigcup_{x \in E} G_x$, hence $E \subset \bigcup_{x \in E} G_x$.
\end{proof}

\begin{definition}[Closed Set]
	$F \subset \mathbb{R}^p$ is called closed if and only if $F^c$ is open.
\end{definition}

\begin{proposition}
	$F \subset \mathbb{R}^p$ is closed if and only if it contains all its boundary points.
\end{proposition}
\begin{proof}
	($\implies$ ) Let $F$ be closed and $x$ is boundary point of $F$. If $x \not\in F$, $x \in F^c$. Hence $x$ is an exterior point. Hence $x \in F$.
	
	($\impliedby$ ) Suppose now $F$ contains all its boundary points. Let $x \not\in F$. Hence $x$ is not interior point of $F$. $x$ is not boundary by assumption. Hence $x$ is exterior. This means $x \in F^c \implies x $ is interior to $F^c$. Hence $F^c$ is open.
\end{proof}

\subsection{Properties of closed sets}

\begin{proposition}
	\quad
	\begin{enumerate}
		\item $R^{p}, \emptyset $ are closed.
		\item Union of a finite family of closed sets is closed.
		\item Intersection of any family of closed sets is closed.
	\end{enumerate}
\end{proposition}
\begin{proof}
	\begin{enumerate}
		\item $(R^p)^c = \emptyset , \emptyset ^c = \mathbb{R}^p$.
		\item $F_1, \ldots, F_N$ are closed. Then \[
		\left( \bigcup_{i = 1} ^N F_i \right) ^c = \bigcap_{i = 1} ^N F_i^c
		.\] 
		Therefore  $ \left( \bigcup_{i = 1} ^N F_i  \right) ^c$ open by property of open sets. Hence $ \bigcup_{i = 1} ^N F_i  $ closed.
	\item Similarly,
		 \[
			 \left( \bigcap_{\alpha \in A} F_\alpha \right) ^c = \bigcup_{\alpha \in A} F_\alpha^c
		.\] 
	\end{enumerate}
\end{proof}

\begin{example}
	\begin{enumerate}
		\item $\{x\} $ is a closed set. 
		\item More generally, any closed ball \[
			\overline{B_r(x)}  = \{y: \|y-x\|\le r\} , r\ge 0
		\] is closed.
		\begin{proof}
			$\overline{B_r(x)}^c$ is the exterior of the ball which we know is open.
		\end{proof}
	\item Cantor set is closed. 
		\begin{proof}
		Each $F_n$ is a finite union of closed union of closed balls, which is closed. $F$ is a countable intersection of closed sets, which is again closed.
		\end{proof}
	\end{enumerate}
\end{example}

